Este trabajo presenta un análisis experimental y comparativo del rendimiento práctico frente a las cotas asintóticas teóricas de cuatro algoritmos de ordenamiento (\textit{MergeSort}, \textit{QuickSort} con partición de Hoare, \textit{PatienceSort} y \texttt{std::sort}) y dos enfoques de multiplicación matricial cuadrada (el método iterativo tradicional \textit{Naive} $\Theta(n^3)$ y el algoritmo \textit{Divide y Vencerás} de Strassen $\mathcal{O}(n^{2.81})$). Se evalúan tiempos de ejecución en CPU de alta precisión y consumos máximos de memoria RAM bajo diversas dimensiones ($N$), topologías de datos y dominios. Los resultados evidencian que en ordenamiento las implementaciones se ajustan estrechamente a $\mathcal{O}(n \log n)$, destacando la baja constante oculta de \texttt{std::sort}. En contraste, en multiplicación matricial la elevada constante oculta y la dispersión en memoria de Strassen provocan que el método \textit{Naive} resulte significativamente superior para todas las instancias evaluadas ($N \le 256$), demostrando el impacto determinante de la jerarquía de memoria física por sobre la complejidad asintótica en rangos prácticos.